\section{Presentations and local finite presentability}
\label{app:presentability}

We give a direct proof of the presentability assertion used in
Lemma~\ref{lem:filtered}. The point requiring care is that an equation
between boundary cells can create new composable pairs. Presentations
must therefore generate defined terms and equations together; taking
an ordinary quotient of the cells already present would not suffice.
All sets and constructions below lie in the fixed universe.

\subsection{Operations and their domains}

Fix $n\in\{1,2,\ldots,\infty\}$, with $\infty$ denoting $\omega$.
Use one sort $C_d$ for each $0\le d\le n$, or for each nonnegative
integer $d$ when $n=\infty$. The total operations are
\[
 s_d,t_d:C_d\longrightarrow C_{d-1}\quad(d>0),
 \qquad
 e_d:C_d\longrightarrow C_{d+1}\quad(d<n).
\]
For $0\le k<d$, write $s_k^d,t_k^d$ for the iterated source and
target, and $e_k^d$ for the iterated identity from sort $k$ to
sort $d$. Each is a finite term. There is a binary partial operation
$*_k^d$ on $C_d$ whose domain is exactly
\begin{equation}\label{eq:pres-domain}
 s_k^d(x)=t_k^d(y).
\end{equation}
Thus $x*_k^d y$ means $x$ after $y$. Superscripts will be omitted
when the sort is clear.

For clarity, the equations imposed are the usual strict globular
category equations. Globularity and the identity boundary equations are
\[
 s_{d-1}s_d=s_{d-1}t_d,\qquad
 t_{d-1}s_d=t_{d-1}t_d,\qquad
 s_{d+1}e_d=\id=t_{d+1}e_d.
\]
The globularity equations apply for $2\le d\le n$, and the identity
boundary equations for $0\le d<n$; when $n=\infty$, there is no upper
bound on these integer indices.
For a pair satisfying \eqref{eq:pres-domain}, the composition boundary
equations are
\[
\begin{array}{ll}
 s_d(x*_{d-1}y)=s_d(y),&
 t_d(x*_{d-1}y)=t_d(x),\\[2pt]
 s_d(x*_k y)=s_d(x)*_k s_d(y),&
 t_d(x*_k y)=t_d(x)*_k t_d(y)\quad(k<d-1).
\end{array}
\]
For each $k$, associativity is imposed on triples with
$s_k(x)=t_k(y)$ and $s_k(y)=t_k(z)$. The unit equations are
\[
 x*_k e_k^d(s_k^d x)=x,\qquad
 e_k^d(t_k^d x)*_k x=x.
\]
Identities preserve every lower composition:
$e_d(x*_k^d y)=e_d(x)*_k^{d+1}e_d(y)$ whenever the pair is
composable and $d<n$.

Interchange, for $k<l<d$, is
\begin{equation}\label{eq:pres-interchange}
 (x*_l y)*_k(z*_l w)=(x*_k z)*_l(y*_k w).
\end{equation}
Its guard includes all four inner matching equations:
\begin{equation}\label{eq:pres-interchange-guard}
\begin{aligned}
 s_l^d(x)&=t_l^d(y),&\qquad s_l^d(z)&=t_l^d(w),\\
 s_k^d(x)&=t_k^d(z),& s_k^d(y)&=t_k^d(w).
\end{aligned}
\end{equation}
These equations also imply the two outer matching conditions.
For the left side, globularity gives
\[
 s_k^d(x*_l y)=s_k^d(x)=s_k^d(y),\qquad
 t_k^d(z*_l w)=t_k^d(z)=t_k^d(w).
\]
For the right side, the iterated boundary formulas give
\[
\begin{split}
 s_l^d(x*_k z)&=s_l^d(x)*_k^l s_l^d(z),\\
 t_l^d(y*_k w)&=t_l^d(y)*_k^l t_l^d(w).
\end{split}
\]
The lower-dimensional composites in these formulas are defined by
the last two equations of \eqref{eq:pres-interchange-guard} and
globularity. The first two equations then make their values equal.
Thus the guard describes an ordinary composable array, not an
additional restriction on interchange.
Definedness of the right-hand side alone is not used to infer
the guard or definedness of the left-hand side.

These are a set of conditional equations with finite terms and
finitely many premises. In the $\omega$ case the sorts and operation
symbols are countable, but no operation or equation involves infinitely
many arguments. The resulting models are precisely strict globular
$n$-categories. For finite $n$, adjoining iterated identities above
dimension $n$ gives the convention used in the body of the paper.

\subsection{Defined terms and presentations}

A \emph{sorted set of generators} $G$ assigns a dimension to each
generator. Let $T(G)$ be the set of all finite well-sorted raw terms
in these generators and the operation symbols. A formal composite
belongs to $T(G)$ even if its boundary equation is not yet known;
it will represent a cell only after being declared defined.

Let $R$ be a set of equations between raw terms of the same sort.
An equation in $R$ is a \emph{strong equation}: its interpretation
requires both terms to be defined and their values to agree. In
particular, $t=t$ can be used to require that a particular term be
defined.

\begin{lemma}[Strict-category presentations]
\label{lem:strict-presentations}
For every sorted set $G$ and every set $R$ of strong equations,
there is a small strict $n$-category $P(G,R)$ such that, naturally in
$Y$, strict functors $P(G,R)\to Y$ correspond to sorted assignments
of the generators to cells of $Y$ satisfying $R$.
\end{lemma}

\begin{proof}
Generate simultaneously judgments $\mathsf D(t)$ and $t=u$,
meaning that $t$ is defined and that $t,u$ are defined with the
same value. Use the following rules.
\begin{enumerate}[label=(\roman*)]
\item Every generator is defined, and each member of $R$ is an
equation.
\item An equation $t=u$ implies $\mathsf D(t)$ and $\mathsf D(u)$.
Definedness of an operation term implies definedness of all its
arguments. Definedness of $x*_k y$ also implies its matching
equation \eqref{eq:pres-domain}.
\item Total operations applied to defined arguments give defined
terms. Defined arguments satisfying \eqref{eq:pres-domain} give a
defined composite.
\item Equality is symmetric and transitive. Reflexivity is the
rule $\mathsf D(t)\Rightarrow t=t$; it is not imposed on undefined
raw terms.
\item Equal arguments transport definedness and equality through
an operation. Precisely, from $t_i=u_i$ for all arguments and
$\mathsf D(f(t_1,\ldots,t_r))$, infer
$\mathsf D(f(u_1,\ldots,u_r))$ and equality of these two operation
terms.
\item Impose every ground instance of the strict category equations
above when its arguments are defined and its stated matching
conditions hold. The conclusion is equality of the two sides and
therefore definedness of both. In particular, interchange retains
all four conditions in \eqref{eq:pres-interchange-guard}.
\end{enumerate}
Globularity, identity boundary and unit equations require just
definedness of their variables. Composition boundary equations and
preservation of composition by identities require
\eqref{eq:pres-domain}; associativity requires its two matching
equations. Thus every rule has finitely many premises. All rules
are sound in a strict category, including the definedness conclusions
in the last rule.

Start with the assumptions and repeatedly add the conclusions of
all applicable rules. The union of the finite stages is the least
closed set of judgments: any rule's finitely many premises occur
together at some stage. In particular, every resulting judgment
has a finite derivation. This construction involves a set of
judgments, since the signature and the set of raw finite terms
are sets.

Let $P(G,R)_d$ be the defined terms of sort $d$ modulo derivable
equality. The equality rules give an equivalence relation on this
set. Total operations descend by rules (iii) and (v). For the
partial operations, suppose the boundary classes of $[x]$ and
$[y]$ match. Total boundary terms are defined, so equality of
their classes means that $s_k(x)=t_k(y)$ is derivable. Rule (iii)
then declares $x*_k y$ defined.

This construction does not depend on the representatives. Indeed,
if $x=x'$ and $y=y'$ are derivable, total-operation congruence gives
\[
 s_k(x)=s_k(x'),\qquad t_k(y)=t_k(y').
\]
The matching equality for $x,y$ therefore gives the matching
equality for $x',y'$. Rule (iii) defines both composites, and
rule (v) equates them. Conversely, rule (ii) ensures that every
admitted composite has matching boundaries. The operation domain
in the quotient is consequently exactly the set of matching
pairs. Rule (vi) gives all the strict category laws there.

Now let an assignment $G\to Y$ satisfy $R$. Induction on derivations
shows that each term declared defined evaluates in $Y$ and each
derived equation has equal values. For rule (v), equal arguments
have identical boundary values, so existence of one partial
operation value entails existence of the other. Evaluation
therefore descends to a strict functor $P(G,R)\to Y$.
It is unique by induction on terms, since definedness is hereditary
to subterms and the images of the generators are fixed. Conversely,
every strict functor gives such an assignment. This proves the
universal property.

The constructed object is small. There is no inconsistency
obstruction: the strict category with one cell in each sort
satisfies every strong equation under the unique assignment.
If $G$ is empty then there are no raw ground terms, since all
operations have positive arity. The construction gives the empty
strict category in that case.
\end{proof}

\begin{remark}\label{rem:presentations-create-composites}
The simultaneous construction is necessary even for ordinary
categories. If the source and target of the free arrow are
identified, its presentation must create the square of the arrow
and all higher powers. These are not obtained by merely identifying
members of the original set of cells. Rule (iii) creates them as
soon as their boundary equalities become derivable.
\end{remark}

\subsection{Limits, colimits, and finite presentations}

\begin{proposition}[Local finite presentability]
\label{prop:strict-presentability}
For every $1\le n\le\infty$, the category $\nCat$ is locally small,
complete, cocomplete, and locally finitely presentable. Filtered
colimits are computed on cells. Every set of objects has a common
infinite regular presentability bound.
\end{proposition}

\begin{proof}
\emph{Colimits.}
For a small diagram $X:I\to\nCat$, use a generator $[i,a]$ for
each cell $a$ of each $X_i$. Impose all its operation-table
equations:
\[
\begin{split}
 s([i,a])&=[i,s(a)],\qquad t([i,a])=[i,t(a)],\\
 e([i,a])&=[i,e(a)],\qquad
 [i,a]*_k[i,b]=[i,a*_k b]
\end{split}
\]
where the last equation is included for every defined composite.
For each arrow $f:i\to j$ also impose
$[i,a]=[j,X(f)(a)]$. All these generators and equations form
sets. By Lemma~\ref{lem:strict-presentations}, maps from this
presentation to $Y$ are exactly compatible families of strict
functors $X_i\to Y$. It is therefore the colimit.

\emph{Limits and local smallness.}
In each dimension, take the set of compatible families in the
diagram. Source, target and identity are coordinatewise.
Two families have matching boundaries exactly when their
coordinates do, so coordinatewise composition is defined on
exactly the required domain and again gives a compatible family.
All laws and the limit universal property hold coordinatewise.
This includes the empty limit, with a singleton in each sort.
Strict functors form a subset of the product of the sets of
cell maps, so the category is locally small.

\emph{Filtered colimits and finite witnesses.}
For a filtered diagram, form the colimit of each set of cells.
Finitely many representatives occur at a common stage. If a
matching equation holds in the colimit, it holds at a later
common stage, where the composite can be formed. A further
common stage proves independence of choices. Each category law
holds at a stage containing its finite tuple and its matching
equations. This constructs the colimit with its universal
property, without any injectivity assumption on transition maps.

We will also use the finite-term version of this argument.
If finitely many assigned generator values make a finite
collection of terms defined and satisfy specified equations in
the colimit, these facts hold at a common stage. Induct on the
finitely many subterms: first lift the arguments, then the
boundary equality needed for a partial composite, and then
the composite. Finally lift the required equalities. Each step
uses the equality criterion for filtered colimits of sets.

\emph{Finite presentations.}
If $G$ and $R$ are finite, a map
$P(G,R)\to\colim_iX_i$ is an assignment of finitely many generator
values satisfying finitely many strong equations. The preceding
finite-term argument lifts that assignment and all the relations
to a single stage. Lemma~\ref{lem:strict-presentations} then
gives a functor to that stage. If two stage functors become
equal in the colimit, their values on the finite set $G$ agree
at a later common stage; their entire functors agree there by
the presentation universal property. Thus
\[
 \colim_i\Hom(P(G,R),X_i)
       \ \cong\ \Hom(P(G,R),\colim_iX_i).
\]
Hence every finite presentation is finitely presentable.
This does not assert that it has finitely many cells.

\emph{A generating set of finite presentations.}
Every strict category $X$ has its full cell presentation:
one generator $g_a$ for each cell and all operation-table
equations. Evaluation identifies its presented category with
$X$. Indeed it is onto, and induction on a defined term,
using the operation-table equations, equates that term to
the generator of its value. This also proves injectivity.

Write this presentation as $(G_X,R_X)$. Consider all finite
pairs $(G_0,R_0)$ with $G_0\subseteq G_X$, $R_0\subseteq R_X$,
and with every generator occurring in $R_0$ contained in
$G_0$. Ordered by inclusion, these pairs form a nonempty
filtered poset: the empty pair is allowed, and finite unions
give upper bounds. Their presented categories have colimit
$P(G_X,R_X)$. To verify this, a compatible family of maps
from the finite presentations assigns a value to every
generator and enforces every relation, since each appears
in some finite pair. The presentation universal property
gives exactly a map from $P(G_X,R_X)$.

The finite pairs need not be closed under all operations.
Finite presentations themselves supply the required closures.
There is a set of finite presentation types: use finite
numbered generator sets with specified sorts and finite sets
of raw term equations. Thus every object is a filtered
colimit of objects from a set of finitely presentable objects.
Together with cocompleteness, this is local finite
presentability. The same reasoning applies to the countable
signature in the $\omega$ case.

\emph{A common regular bound.}
More generally, if $G$ and $R$ have cardinality less than an
infinite regular cardinal $\kappa$, the same proof makes
$P(G,R)$ $\kappa$-presentable. In a $\kappa$-filtered diagram
one lifts fewer than $\kappa$ generator values and the
finitely many subterms of each of fewer than $\kappa$
relations. All stage, morphism and equality witnesses involved
form a diagram of size less than $\kappa$, so a cocone places
them at one stage. Equality of maps is detected by the
fewer than $\kappa$ generator equalities in the same way.

The full cell presentation of $X$ has cardinality at most
\[
 \max\left\{\aleph_0,\left|\coprod_d X_d\right|\right\},
\]
because the signature is countable and operations have arity
at most two. Choose a larger infinite regular cardinal.
The preceding argument makes $X$ presentable at that bound.
For a set of objects, bound the cardinalities of all their
full presentations first and choose one larger infinite
regular cardinal. This gives the claimed common bound.
\end{proof}
